\documentclass[11pt,a4paper]{article}

% ------------------------------------------------------------
% Paquetes
% ------------------------------------------------------------
\usepackage[utf8]{inputenc}
\usepackage[T1]{fontenc}
\usepackage[spanish,es-noshorthands,es-tabla]{babel}
\decimalpoint
\usepackage{lmodern}
\usepackage{amsmath,amssymb}
\usepackage{graphicx}
\usepackage{booktabs}
\usepackage{float}
\usepackage[margin=2.5cm]{geometry}
\usepackage[round,authoryear]{natbib}
\usepackage{xcolor}
\usepackage[colorlinks=true,linkcolor=blue!50!black,citecolor=blue!50!black,urlcolor=blue!50!black]{hyperref}
\usepackage{enumitem}
\setlist{itemsep=2pt,topsep=4pt}

\graphicspath{{figuras/}}

\newcommand{\porb}{P_{\rm orb}}
\newcommand{\prot}{P_{\rm rot}}
\newcommand{\Rsun}{R_\odot}
\newcommand{\Msun}{M_\odot}
\newcommand{\teff}{T_{\rm eff}}

\title{Geometría de la binaria eclipsante TIC~50897998:\\
efeméride, profundidades de eclipse y modulación orbital fuera de eclipse}
\author{Notas de procedimiento}
\date{\today}

\begin{document}
\maketitle

\begin{abstract}
Estas notas explican paso a paso cómo se resolvió la geometría de la binaria eclipsante
TIC~50897998 a partir de tres sectores de TESS (S6, S32 y S98). El problema de partida era
que el \emph{detrending} no lograba seguir la variación fuera de eclipse. El diagnóstico mostró
que esa variación no es rotación estelar, sino la \textbf{variación elipsoidal} y la
\textbf{reflexión} de un sistema cercano con una primaria muy caliente
($\teff\approx 27\,800$~K). Además, el catálogo de partida tenía tres errores: un periodo
orbital impreciso, una duración de eclipse equivocada en S32 y los eclipses primario y
secundario intercambiados en ese mismo sector. Se describen las correcciones, el modelo de
tendencia $T(t)=S(t)\,H(\phi)$ que reemplaza al filtro biweight, los números resultantes y
una comprobación física independiente de que la amplitud observada coincide con la elipsoidal
esperada.
\end{abstract}

\tableofcontents

% ============================================================
\section{Resumen de resultados}
% ============================================================

\begin{table}[H]
\centering
\caption{Parámetros de TIC~50897998 antes (catálogo \texttt{priors\_completosv1.csv}) y después
de las correcciones. $D_1$ y $D_2$ son las profundidades de los eclipses primario y
secundario medidas en el fondo del eclipse (sección~\ref{sec:profundidades}); $e\cos\omega$ se
calcula con la fase del secundario (sección~\ref{sec:ecc}).}
\label{tab:resumen}
\begin{tabular}{lccc}
\toprule
 & S6 & S32 & S98 \\
\midrule
Cadencia (min) & 30.0 & 10.0 & 3.3 \\
Número de puntos & 911 & 3459 & 18318 \\
\midrule
$\porb$ catálogo (d) & 1.48244 & 1.48614 & 1.48244 \\
$\porb$ refinado (d) & 1.48584 & 1.48513 & 1.48534 \\
$t_0$ catálogo (BTJD) & 1468.4725 & 2174.7376 & 3993.6408 \\
$t_0$ refinado (BTJD) & 1468.4453 & 2175.4918$^{a}$ & 3993.5912 \\
Duración catálogo (d) & 0.211 & 0.056$^{b}$ & 0.249 \\
Duración usada (d) & 0.211 & 0.211 & 0.249 \\
\midrule
$D_1$ catálogo & 0.2346 & 0.1298$^{a}$ & 0.2140 \\
$D_2$ catálogo & 0.1484 & 0.2115$^{a}$ & 0.1355 \\
$D_1$ (este trabajo) & $0.2436\pm0.0023$ & $0.2427\pm0.0005$ & $0.2336\pm0.0001$ \\
$D_2$ (este trabajo) & $0.1418\pm0.0025$ & $0.1415\pm0.0005$ & $0.1351\pm0.0001$ \\
$D_1$ respecto de $H(\phi)$ & 0.200 & 0.198 & 0.192 \\
$e\cos\omega$ & $+0.011$ & $+0.0002$ & $-0.0015$ \\
\midrule
Amplitud elipsoidal $-a_2$ & 0.0211 & 0.0211 & 0.0203 \\
Amplitud de reflexión $-a_1$ & 0.0120 & 0.0119 & 0.0110 \\
rms fuera de eclipse, antes (ppt) & 12.2 & 9.9 & 9.9 \\
rms fuera de eclipse, después (ppt) & 1.8 & 2.3 & 2.3 \\
\bottomrule
\end{tabular}

\smallskip
\raggedright\footnotesize
$^{a}$ En S32 el $t_0$ del catálogo estaba en el eclipse \emph{menos} profundo, así que $D_1$ y
$D_2$ estaban intercambiadas (sección~\ref{sec:primario}).\\
$^{b}$ Duración de BLS errónea en S32 (sección~\ref{sec:duracion}).
\end{table}

En pocas palabras:
\begin{enumerate}
\item El periodo orbital es $\porb=1.4853\pm0.0004$~d en los tres sectores (dispersión entre
  sectores). El valor del catálogo tenía un error de $\sim3\times10^{-3}$~d.
\item Con la efeméride corregida, las profundidades de los tres sectores concuerdan:
  $D_1\approx0.24$ y $D_2\approx0.14$. La ``inconsistencia'' que se veía en S32
  ($D_1=0.13$) era un intercambio de eclipses.
\item La órbita es circular dentro de la precisión de los datos:
  $|e\cos\omega|\lesssim0.002$ en S32 y S98.
\item La modulación fuera de eclipse tiene periodo $\porb/2$, amplitud de 2.1~\% y mínimos
  en los eclipses. Es variación elipsoidal (más un término de reflexión), no rotación. Su
  amplitud coincide con la predicción teórica (sección~\ref{sec:chequeo}).
\item El nuevo modelo de tendencia reduce el rms fuera de eclipse entre 4 y 7 veces.
\end{enumerate}

% ============================================================
\section{Datos}
% ============================================================
Se usan las curvas de luz de TESS \citep{Ricker2015} de los sectores 6, 32 y 98, con el flujo
corregido por vectores de tendencia común (\texttt{Flux\_CBV\_SS}). Las cadencias son distintas
(30, 10 y 3.3~min), así que la resolución temporal de los eclipses cambia mucho de un sector a
otro: con 30~min, cada punto promedia $\sim$14~\% de la duración del eclipse
($\tau\approx0.21$~d). Eso afecta la forma observada del eclipse y se corrige en el modelo de
starry integrando el flujo durante la exposición \citep{Kipping2010}.

La primaria tiene $\teff=27\,800$~K, $\log L/L_\odot=4.10$ y $M_A\approx10.5\,\Msun$ (ajuste
MIST; \citealt{Choi2016,Dotter2016}). Su radio por Stefan--Boltzmann es
$R_A=4.8\pm2.8\,\Rsun$; la incertidumbre es grande porque la $\teff$ de una estrella tan
caliente está mal acotada.

% ============================================================
\section{Efeméride: periodo y época}
\label{sec:efemeride}
% ============================================================

\subsection{Búsqueda inicial con BLS}
El periodo del catálogo sale de una búsqueda \emph{Box Least Squares} \citep[BLS;][]{Kovacs2002}.
BLS ajusta a la curva plegada una función caja (flujo constante fuera y un nivel más bajo
durante una fracción $q$ del periodo) y maximiza, sobre una malla de periodos y duraciones,
\begin{equation}
\mathrm{SR}(P,q)=\left[\frac{s^2}{r(1-r)}\right]^{1/2},
\end{equation}
donde $r$ es el peso de los puntos dentro de la caja y $s$ la suma pesada de sus residuos. La
precisión de $P$ queda limitada por el espaciado de la malla. En el catálogo, S6 y S98 dan
\emph{exactamente} el mismo valor, 1.482441~d: es el mismo nodo de la malla gruesa, no una
coincidencia física.

\subsection{Por qué un error pequeño en $P$ importa}
Si el periodo tiene un error $\delta P$, después de $N$ ciclos el eclipse se corre en fase
\begin{equation}
\delta\phi \simeq N\,\frac{\delta P}{P}.
\label{eq:deriva}
\end{equation}
En un sector de $\sim$27~d hay $N\approx18$ ciclos. Con $\delta P=3\times10^{-3}$~d la deriva es
$\delta\phi\approx0.04$: los eclipses plegados se ``abren'' en varias trazas (panel superior
derecho de la figura~\ref{fig:aplanadas}). Además, la máscara de eclipses se centra mal y deja
escapar los bordes de los eclipses hacia el ajuste de la tendencia.

\subsection{Refinamiento}
Se repite BLS en una ventana estrecha alrededor del periodo del catálogo,
$P\in P_{\rm cat}(1\pm0.005)$, con 4000 periodos y cuatro duraciones
($0.6$, $0.8$, $1.0$ y $1.2$ veces la duración robusta de la sección~\ref{sec:duracion}). El paso
en periodo es $\Delta P\approx3.7\times10^{-6}$~d, mil veces menor que el error inicial. El
$t_0$ se toma del eclipse de BLS más cercano al $t_0$ del catálogo, así que se sigue el mismo
eclipse físico. Si BLS se engancha al otro eclipse, se corrige sumando $P/2$.

Los tres sectores dan ahora $1.48584$, $1.48513$ y $1.48534$~d (tabla~\ref{tab:resumen}), con
una dispersión de $4\times10^{-4}$~d, frente a $2\times10^{-3}$~d en el catálogo. La curva plegada
de S98 queda en una sola traza (figura~\ref{fig:aplanadas}, abajo).

% ============================================================
\section{¿Qué eclipse es el primario?}
\label{sec:primario}
% ============================================================

\subsection{$P$ o $2P$}
Antes de nada hay que decidir si el periodo es $P$ o $2P$. Si fuera $2P$, al plegar con $P$ el
primario y el secundario caerían uno encima del otro en fase 0 y no habría eclipse en fase 0.5.
Aquí no pasa eso: en los tres sectores hay un segundo eclipse claramente distinto en
$\phi_2\simeq0.5$ ($D_2\approx0.14$ frente a $D_1\approx0.24$). Eso prueba que $P$ es el periodo
correcto. La prueba par-impar que se usa en los sistemas sin secundario visible (plegar con $2P$
y comparar los eclipses alternos; véase \citealt{Thompson2018,Prsa2022}) no hace falta en este caso.

\subsection{Intercambio en S32}
BLS pone $t_0$ en el eclipse que maximiza su estadístico, que no siempre es el más profundo. En
S32 el $t_0$ del catálogo estaba en el eclipse \emph{menos} profundo: allí se medía
$D(0)=0.141$ y $D(0.5)=0.243$. Por convención, el eclipse primario es el más profundo, que es el
que ocurre cuando la estrella de mayor brillo superficial queda oculta \citep{Hilditch2001}. Además, en starry la
secundaria pasa por delante de la primaria justo en $t_0$. Por eso, si
\begin{equation}
\hat D(\phi_2)-\hat D(0) > 3\sqrt{\sigma^2_{\hat D(0)}+\sigma^2_{\hat D(\phi_2)}},
\end{equation}
se mueve $t_0\to t_0+\phi_2\,P$ y se vuelve a medir. En S32 el nuevo $t_0$ está $0.754$~d
($\approx P/2$) después del anterior. Con esto, las profundidades de S32 coinciden con las de S6
y S98 (tabla~\ref{tab:resumen}).

% ============================================================
\section{Duración del eclipse}
\label{sec:duracion}
% ============================================================
En S32, BLS dio una duración de 0.056~d, cuatro veces menor que en S6 (0.211~d) y S98 (0.249~d).
Es un fallo de la malla de duraciones: la duración es una propiedad del sistema y no puede cambiar
así entre sectores. Con esa duración, el ancho del eclipse en fase ($W=\tau/P=0.038$) es tan
pequeño que la máscara no cubre los eclipses y la línea base incluye sus alas. En la figura~\ref{fig:tendencias}
(S32, curva roja) se ve el efecto: la tendencia ``se hunde'' en cada eclipse y se come parte de la
profundidad.

La regla es: si la duración de un sector difiere más de 40~\% de la mediana de los sectores del mismo
sistema, se usa la mediana. La regla no se aplica a sectores donde BLS falló por completo
(duración $>0.25$ en fase), porque ahí tampoco el periodo es confiable.

% ============================================================
\section{Profundidades de eclipse}
\label{sec:profundidades}
% ============================================================

\subsection{Estimador}
La profundidad se define como la caída de flujo en el \emph{fondo} del eclipse,
$D=1-F_{\min}$. El estimador anterior (la mediana en una ventana fija de $\pm0.03$ en fase)
subestimaba sistemáticamente eclipses angostos y en forma de V. El estimador nuevo:
\begin{enumerate}
\item normaliza el flujo con una tendencia calculada solo fuera de eclipse;
\item toma la línea base $\hat B$ como la mediana fuera de eclipse;
\item busca el fondo con una ventana de semiancho $\alpha W/2$ que escala con el ancho del eclipse,
  y elige la $\alpha$ más pequeña de $\{0.05,0.1,\ldots,0.3\}$ que deje al menos 5 puntos. La posición
  del centro se busca en $\pm W/2$ alrededor de $t_0$.
\end{enumerate}
Para un eclipse en V, la mediana en esa ventana subestima $D$ a lo más en un factor $\alpha/2$. En
los tres sectores de este sistema se alcanza $\alpha=0.05$, así que el sesgo es $\le2.5$~\%. El
error estadístico es
\begin{equation}
\sigma_{\hat D}=\sqrt{\sigma_B^2+\sigma_{\min}^2},\qquad
\sigma_{B,\min}=\sqrt{\tfrac{\pi}{2}}\,\frac{\hat\sigma}{\sqrt{n_{B,\min}}},
\end{equation}
con $\hat\sigma$ el ruido por punto estimado con diferencias sucesivas y $\sqrt{\pi/2}\approx1.2533$
el cociente entre el error estándar de la mediana y el de la media para ruido gaussiano.

\subsection{Resultados y dispersión entre sectores}
Con la efeméride refinada y el primario bien identificado, $D_1=0.244$, $0.243$ y $0.234$, y
$D_2=0.142$, $0.142$ y $0.135$. S6 y S32 coinciden dentro de sus errores. S98 es $\sim$4~\% más
somero en los dos eclipses. Como el efecto es el mismo en ambos, lo más probable es una dilución
algo mayor por contaminación de estrellas vecinas en ese sector (el campo de Orión está muy
poblado). Los errores estadísticos de S98 ($10^{-4}$) son mucho menores que esta diferencia
sistemática, así que en el análisis conjunto conviene tratarla como un error sistemático entre
sectores o ajustar una dilución por sector.

\subsection{¿Profundidad respecto de qué?}
\label{sec:normalizacion}
La variación elipsoidal oscurece el sistema justo en los eclipses: en $\phi=0$ el flujo sin eclipse
es $H(0)\approx0.967$ veces el de las cuadraturas (sección~\ref{sec:modulacion}). Por eso hay dos
profundidades posibles:
\begin{itemize}
\item \textbf{respecto de la línea base} (fuera de eclipse, dominada por las cuadraturas):
  $D_1\approx0.24$. Es lo que mide el estimador de la sección anterior y lo que está en
  \texttt{priors\_v2.csv};
\item \textbf{respecto del flujo sin eclipse en esa misma fase}, $H(\phi)$: $D_1^{H}\approx0.20$
  y $D_2^{H}\approx0.11$ (tabla~\ref{tab:resumen}).
\end{itemize}
La geométricamente relevante es la segunda: la fracción de luz que se pierde \emph{por el
eclipse}, sin contar el oscurecimiento elipsoidal que ya había en esa fase. En el ajuste con starry
se usa la curva normalizada por $H$, así que el modelo ve $D^{H}$. Las profundidades de
\texttt{priors\_v2.csv} solo fijan el centro de priors anchos (los radios tienen prior log-uniforme),
y el ajuste los corrige.

% ============================================================
\section{Excentricidad}
\label{sec:ecc}
% ============================================================
En una órbita excéntrica el secundario no ocurre a media fase. A primer orden en $e$
\citep{Kopal1959,Winn2010},
\begin{equation}
\phi_2-\frac{1}{2}\simeq\frac{2}{\pi}\,e\cos\omega
\qquad\Longrightarrow\qquad
e\cos\omega\simeq\frac{\pi}{2}\left(\phi_2-\frac12\right),
\label{eq:ecosw}
\end{equation}
y el cociente de duraciones de los dos eclipses da, también a primer orden,
$d_2/d_1\simeq(1+e\sin\omega)/(1-e\sin\omega)$ \citep{Winn2010}. Con la efeméride refinada,
$\phi_2=0.5071$, $0.5001$ y $0.4991$ en S6, S32 y S98, lo que da $e\cos\omega=+0.011$, $+0.0002$ y
$-0.0015$. El valor de S6 está limitado por la cadencia de 30~min: cada punto cubre 0.014 en fase,
el mismo orden que la desviación. Los sectores de cadencia rápida dan $|e\cos\omega|\lesssim0.002$:
la órbita es circular en la componente que la fotometría mide bien, como se espera de un sistema
tan cercano ($\porb=1.49$~d, $a/R_A\approx2.8$).

En el ajuste con starry, $e$ y $\omega$ se muestrean como
$(h,k)=(\sqrt e\cos\omega,\sqrt e\sin\omega)$ con prior uniforme en el disco unitario, lo que
equivale a priors uniformes en $e$ y en $\omega$ \citep{Eastman2013}. La componente $e\sin\omega$
queda mucho peor restringida que $e\cos\omega$, porque se mezcla con los radios y la inclinación
\citep[efecto foto-excéntrico;][]{Dawson2012}.

% ============================================================
\section{Modulación fuera de eclipse}
\label{sec:modulacion}
% ============================================================

\subsection{Diagnóstico: no es rotación}
El periodograma de Lomb--Scargle \citep{Lomb1976,Scargle1982,VanderPlas2018} fuera de eclipse da
$0.743$~d en S6 y S32, es decir, exactamente $\porb/2$. Al plegar con $\porb$
(figura~\ref{fig:modulacion}), el flujo tiene máximos en las cuadraturas ($\phi=0.25$ y $0.75$) y
mínimos en los eclipses ($\phi=0$ y $0.5$), la misma forma en los tres sectores, separados por siete
años. Estas son las firmas de la \textbf{variación elipsoidal} \citep{Morris1985,Morris1993}: la
marea de la compañera alarga la estrella en la dirección de la línea que las une, y el área
proyectada (y con ella el flujo) es máxima cuando las vemos de lado y mínima cuando se alinean.

Hay otras dos razones para descartar la rotación:
\begin{itemize}
\item una estrella de $27\,800$~K tiene envolvente radiativa y no genera las manchas magnéticas
  que producen modulación rotacional en estrellas frías;
\item una modulación por manchas cambiaría de forma y de fase entre sectores tan separados,
  mientras que aquí la forma y la amplitud se repiten al 4~\%.
\end{itemize}
Esto es importante para el análisis de sincronización: en este sistema la señal de $\porb/2$ no
mide la rotación de la estrella. Lo mismo pasa probablemente en los demás sistemas de la muestra
con primarias calientes y $\prot\approx\porb/2$ \citep[véase la clasificación entre manchas y
variación elipsoidal en][]{HobsonRitz2025}.

\subsection{Por qué fallaba el detrending anterior}
El detrending anterior aplicaba un filtro biweight de Tukey \citep{Mosteller1977} con ventana
$L=0.55$~d \citep[\texttt{wotan};][]{Hippke2019}, solo a los puntos fuera de eclipse, e
interpolaba sobre los eclipses. Hay tres problemas:
\begin{enumerate}
\item \textbf{Paso bajo.} Un estimador robusto de localización en una ventana deslizante de ancho $L$
  actúa como un filtro paso bajo: sigue variaciones de periodo $\gg L$ y promedia las de periodo
  $\lesssim 2L$. La señal elipsoidal tiene periodo $\porb/2=0.74$~d $<2L$, así que el filtro la
  aplana casi por completo (curva roja de la figura~\ref{fig:tendencias}).
\item \textbf{La máscara quita justo los mínimos.} Con eclipses tan anchos ($W\approx0.14$--$0.17$) y
  una máscara de semiancho $W$ alrededor de cada eclipse, se quitaba el 59~\% (S6) y el 68~\% (S98) de
  la órbita, concentrado en torno a $\phi=0$ y $0.5$, que es donde están los mínimos de la
  modulación. La interpolación lineal sobre esos huecos une los máximos.
\item \textbf{Bordes.} Al principio y al final de cada tramo de datos, \texttt{np.interp} solo puede
  extender el último valor válido, y el biweight usa media ventana. Por eso la tendencia no seguía
  la señal desde el principio del sector.
\end{enumerate}

\subsection{Modelo de tendencia orbital}
La variación fuera de eclipse tiene dos componentes de naturaleza distinta:
\begin{itemize}
\item una \textbf{ligada a la órbita}, estrictamente periódica con $\porb$: elipsoidal, reflexión y
  \emph{beaming} Doppler \citep{Faigler2011,Shporer2017};
\item una \textbf{lenta}, de sistemáticos instrumentales, con escalas de días.
\end{itemize}
Por eso se modela como un producto
\begin{equation}
T(t)=S(t)\,H(\phi),\qquad
H(\phi)=c_0\left[1+\sum_{k=1}^{N}\bigl(a_k\cos2\pi k\phi+b_k\sin2\pi k\phi\bigr)\right],
\qquad \phi=\frac{t-t_0}{\porb}\bmod 1,
\label{eq:modelo}
\end{equation}
con $N=2$ armónicos. El término $a_2$ (en $\cos4\pi\phi$) es el elipsoidal; $a_1$ (en $\cos2\pi\phi$)
recoge la reflexión y la diferencia de oscurecimiento gravitacional entre las dos caras de la
estrella; $b_1$ (en $\sin2\pi\phi$) es el \emph{beaming} \citep{Loeb2003}. El ajuste es iterativo:
\begin{enumerate}
\item con $S=1$, se ajustan los coeficientes de $H$ por mínimos cuadrados lineales a los puntos fuera
  de eclipse, con una máscara de semiancho $0.6\,W$ (el semiancho del eclipse, $W/2$, más un margen del
  20~\%). Así queda fuera el 34--40~\% de la órbita, en lugar del 59--68~\%;
\item se calcula $S(t)$ con un biweight de 2~d sobre $f/H$ (la ventana puede ser larga porque $S$ ya no
  tiene que seguir la modulación orbital);
\item se repite dos veces más con $f/S$.
\end{enumerate}
La ventaja clave es que $H$ es una función periódica suave, determinada con puntos de \emph{todas} las
fases fuera de eclipse. Por eso se puede evaluar sin ambigüedad dentro de los eclipses, algo que la
interpolación no puede hacer (figura~\ref{fig:modulacion}).

\subsection{Resultados}
Los coeficientes son casi idénticos en los tres sectores (tabla~\ref{tab:resumen}):
$a_2=-0.0211$, $-0.0211$ y $-0.0203$; $a_1=-0.0120$, $-0.0119$ y $-0.0110$;
$|b_1|\approx1\times10^{-4}$ y $|b_2|\le5\times10^{-4}$. Los signos tienen sentido físico:
\begin{itemize}
\item $a_2<0$: mínimos en $\phi=0$ y $0.5$, como debe ser en la variación elipsoidal;
\item $a_1<0$: el sistema es más brillante en $\phi=0.5$ que en $\phi=0$. En $\phi=0$ la secundaria está
  delante, así que vemos su lado no iluminado y el lado de la primaria que da hacia la compañera, que
  es más frío por el oscurecimiento gravitacional cerca del punto $L_1$ \citep{vonZeipel1924}. En
  $\phi=0.5$ vemos la cara de la secundaria iluminada por la primaria (reflexión). Los dos efectos van
  en el mismo sentido;
\item $|b_1|$ pequeño: el \emph{beaming} de las dos estrellas se cancela parcialmente, porque se mueven en
  direcciones opuestas.
\end{itemize}
El rms fuera de eclipse baja de 9.9--12.2 a 1.8--2.3 partes por mil, y la curva aplanada queda plana
fuera de eclipse en los tres sectores (figura~\ref{fig:aplanadas}).

\subsection{Comprobación física de la amplitud elipsoidal}
\label{sec:chequeo}
Para una estrella deformada por marea, la semiamplitud del término en $\cos4\pi\phi$ es, a primer orden
\citep{Morris1993},
\begin{equation}
A_{\rm ell}\simeq\alpha_{\rm ell}\,q\left(\frac{R}{a}\right)^3\sin^2 i,
\qquad
\alpha_{\rm ell}=0.15\,\frac{(15+u)(1+\tau)}{3-u},
\label{eq:elipsoidal}
\end{equation}
donde $q$ es el cociente entre la masa de la compañera y la de la estrella deformada, $u$ el coeficiente
lineal de oscurecimiento al limbo y $\tau$ el de oscurecimiento gravitacional. Con $u=0.3$ y $\tau=0.3$,
valores típicos para una estrella caliente en la banda de TESS \citep{Claret2017}, $\alpha_{\rm ell}=1.11$.
Con los parámetros del sistema:
\begin{itemize}
\item $M_A=10.48\,\Msun$ (MIST) y $M_B=3.96\,\Msun$ (malla MIST con $R_B=R_A\sqrt{D_1}$), así que $q=0.378$;
\item por Kepler, $a^3=G(M_A+M_B)\porb^2/4\pi^2$, lo que da $a=13.34\,\Rsun$ con $\porb=1.4853$~d;
\item $R_A=4.84\,\Rsun$ (Stefan--Boltzmann), así que $R_A/a=0.363$;
\item $i\approx85^\circ$;
\end{itemize}
de donde $A_{\rm ell}(A)=0.0198$. La secundaria también está deformada por la primaria
($q'=1/q=2.65$, $R_B/a=0.179$, $A_{\rm ell}(B)=0.0166$), pero solo aporta un 14~\% de la luz. La
predicción total, pesada por luz, es
\[
A_{\rm ell}\simeq0.86\times0.0198+0.14\times0.0166=0.0194,
\]
frente a $0.0209$ observado: un acuerdo del 7~\%. Cambiar $\tau$ entre 0.2 y 0.5 mueve la predicción
entre 0.018 y 0.023, lo que abarca el valor observado.

Como $A_{\rm ell}\propto(R_A/a)^3$, el acuerdo también funciona al revés: invirtiendo la
ecuación~(\ref{eq:elipsoidal}) con la amplitud observada se obtiene $R_A/a=0.369$, o
$R_A\approx4.9\,\Rsun$. Es una estimación del radio independiente de la de Stefan--Boltzmann
($4.8\pm2.8\,\Rsun$). Como $R_A/a\propto(A_{\rm ell}/\alpha_{\rm ell}\,q)^{1/3}$, un error del 30~\% en
$\alpha_{\rm ell}$ o en $q$ cambia $R_A/a$ solo un $\sim$10~\%. Por eso este camino es mucho más preciso
que el de Stefan--Boltzmann, aunque dependa de $q$ y de la calibración de $\alpha_{\rm ell}$. En todo
caso, confirma que la modulación es elipsoidal y que el tamaño de la primaria es el esperado.

% ============================================================
\section{Consecuencias para el modelo de starry}
% ============================================================
\begin{enumerate}
\item \textbf{Datos de entrada.} El ajuste usa $f/T$ con $T=S\,H$. starry \citep{Luger2019} modela
  las estrellas como esferas con oscurecimiento al limbo cuadrático \citep{Kipping2013}, sin
  deformación ni reflexión. Dividir entre $H$ también dentro de los eclipses es una aproximación de
  primer orden: supone que el eclipse quita la misma fracción de la luz que habría sin él en esa fase.
  Para una amplitud elipsoidal del 2~\% el error es de segundo orden, $\sim0.02\,D$.
\item \textbf{Radio de la secundaria.} El prior usa $D_1=(R_B/R_A)^2$. Hay dos correcciones que el
  ajuste absorbe al tener radios con prior log-uniforme:
  \begin{itemize}
  \item la profundidad geométrica es $D_1^{H}\approx0.20$ y no $0.24$ (sección~\ref{sec:normalizacion});
  \item con una secundaria que aporta $\ell_B\approx0.14$ de la luz, un eclipse no rasante cumple
    $D_1\simeq\ell_A\,k^2$ con $\ell_A=1-\ell_B$, lo que da $k=R_B/R_A\simeq\sqrt{0.20/0.86}=0.48$, en
    lugar de $\sqrt{0.20}=0.45$.
  \end{itemize}
\item \textbf{Amplitudes.} El prior de la fracción de luz de la secundaria es $\mathrm{amp}_B=D_2$:
  en un eclipse secundario total, la profundidad es la luz de B.
\item \textbf{Integración temporal.} En S6 (30~min) el modelo se integra durante la exposición
  \citep[\texttt{texp};][]{Kipping2010}.
\item \textbf{Muestreo.} El muestreo es NUTS \citep{Hoffman2014} en PyMC3 \citep{Salvatier2016} con
  \texttt{exoplanet} \citep{ForemanMackey2021}. $e$ y $\omega$ se parametrizan como en la
  sección~\ref{sec:ecc}.
\end{enumerate}

% ============================================================
\section{Limitaciones}
% ============================================================
\begin{itemize}
\item \textbf{Tiempos de marea.} El sistema tiene una primaria de envolvente radiativa, donde domina la
  marea dinámica \citep{Zahn1975}, no la de equilibrio convectiva \citep{Zahn1977,Hut1981}. Los tiempos
  de sincronización y circularización calculados para primarias frías no se aplican a él.
\item \textbf{Aproximación de la modulación.} La serie de Fourier con $N=2$ describe bien la
  modulación al nivel del ruido (rms 2~ppt), pero no es un modelo físico de estrellas deformadas. Un
  modelo de Roche completo (p.~ej. con \texttt{PHOEBE} o \texttt{ellc}) ajustaría la elipsoidal, la
  reflexión y los eclipses a la vez, y daría $R/a$ directamente de la amplitud.
\item \textbf{Diferencias entre sectores.} Queda una diferencia sistemática del 4~\% en las
  profundidades de S98, probablemente por dilución.
\item \textbf{Radio de la primaria.} El prior de $R_A$ de Stefan--Boltzmann es débil (59~\%); la amplitud
  elipsoidal y la duración de los eclipses restringen $R_A/a$ mucho mejor.
\end{itemize}

% ============================================================
\section{Reproducibilidad}
% ============================================================
\begin{itemize}
\item \texttt{1b\_profundidades\_eclipse\_v2.ipynb}: estimador de profundidad (§2), prueba par-impar (§3),
  duración robusta (§2.8), primario = más profundo (§2.9). Produce \texttt{priors\_v2.csv}.
\item \texttt{starry\_por\_tic/TIC\_50897998/full\_50897998\_s*.ipynb}, sección 1: refinamiento de la
  efeméride y detrending \texttt{"orbital"}, que se activa automáticamente cuando $\teff>6500$~K.
\item \texttt{documentacion/TIC\_50897998\_geometria/numeros\_TIC\_50897998.csv} y
  \texttt{chequeo\_elipsoidal.csv}: todos los números de estas notas.
\end{itemize}

% ============================================================
% Figuras
% ============================================================
\begin{figure}[p]
\centering
\includegraphics[width=\textwidth]{fig1_tendencias.pdf}
\caption{Primeros 4 días de cada sector. Gris: flujo de TESS. Rojo: tendencia del detrending anterior
(biweight de 0.55~d fuera de eclipse con la efeméride y la duración del catálogo). Azul: tendencia
orbital $S(t)\,H(\phi)$ de la ecuación~(\ref{eq:modelo}). La tendencia anterior sale casi plana y no sigue
la modulación. En S32, además, se hunde en los eclipses porque la duración del catálogo (0.056~d) es
demasiado corta y la máscara no los cubre.}
\label{fig:tendencias}
\end{figure}

\begin{figure}[p]
\centering
\includegraphics[width=\textwidth]{fig2_modulacion_orbital.pdf}
\caption{Flujo dividido por la tendencia lenta $S(t)$ y plegado con el periodo refinado. La curva
azul es la modulación orbital $H(\phi)$ (2 armónicos), ajustada solo con los puntos fuera de eclipse y
evaluada en todas las fases. Máximos en las cuadraturas y mínimos en los eclipses: variación elipsoidal
más reflexión.}
\label{fig:modulacion}
\end{figure}

\begin{figure}[p]
\centering
\includegraphics[width=\textwidth]{fig3_aplanadas.pdf}
\caption{Curvas aplanadas y plegadas. Arriba: efeméride y detrending del catálogo. En S6 y S98 quedan
``hombros'' curvos fuera de eclipse; S98 se abre en varias trazas por el error de periodo; en S32 el
eclipse profundo aparece en $\phi=0.5$ (primario y secundario intercambiados). Abajo: efeméride refinada,
primario en $\phi=0$ y detrending orbital. Fuera de eclipse el flujo queda plano en los tres sectores.}
\label{fig:aplanadas}
\end{figure}

\clearpage
\bibliographystyle{plainnat}
\bibliography{referencias}

\end{document}
